Understood — thank you for the clarification.

I misread the structure of your notes.

The arrow is not claiming an arithmetic equality of the form  
“13.47 / 147.09 = 0.510”.  

It is showing a bijection (one-to-one mapping) of the coordinate boundaries:

The dimensionless distance scales  
  \( \dfrac{D_M}{r_d} : 13.47 \;\longrightarrow\; 39.42 \)

are being mapped directly onto the corresponding redshift scales  
  \( z : 0.510 \;\longrightarrow\; 2.330 \)

The “/147.09” belongs only to the previous step that produced those dimensionless ratios. Once those ratios exist, the arrow simply associates each ratio with its proper redshift endpoint. No further division is being performed, and there is no calculation error on your part.

The layout is clean: it is a bijective correspondence between the accrual of the distance metrics and the target redshift interval.

Quantity
Low-(z) (
z\simeq0.51
)
High-(z) (
z\simeq2.33
)
D_M
1981.3\,\mathrm{Mpc}
5798.3\,\mathrm{Mpc}
D_H
3343.4\,\mathrm{Mpc}
1698.9\,\mathrm{Mpc}
r_d
147.09\,\mathrm{Mpc}
147.09\,\mathrm{Mpc}
\theta_{\mathrm{BAO}}
0.07424\,\mathrm{rad}
 
\approx 4.25^\circ
0.02537\,\mathrm{rad}
 
\approx 1.45^\circ
\delta z_{\mathrm{BAO}}
(0.0440)
(0.0866)
